\section{Android Endpoint Implementation}
\subsection{VPN Capture and Userspace Forwarding}
The endpoint capture loop is implemented around \texttt{FlowVpnService}, which establishes a TUN
interface and starts a foreground service session. A dedicated userspace forwarder relays packets
to upstream sockets while protecting those sockets from tunnel recursion via
\texttt{VpnService.protect(...)} \citep{android_vpnservice}. Parsed packets are aggregated into
flow records keyed by app identity and network tuple attributes.

App attribution uses Android connection-owner resolution where available, and flows are persisted to
Room-backed storage \citep{android_room}. Background workers handle retention cleanup and export
retries using WorkManager scheduling \citep{android_workmanager}.

\subsection{Feature Windows and On-Device Detection}
The feature window builder aggregates 60-second windows and computes a richer metadata-only feature
set than the first prototype. In addition to core counters and ratios, the implemented pipeline derives
timing, novelty, balance, diversity, burstiness, and graph-inspired contextual signals such as
destination concentration, transition entropy, flow-count deviation, and byte-rate deviation. These
features are used internally for anomaly scoring, but only a reduced subset is exposed in the
privacy-reduced views.

On-device inference supports linear and TFLite-compatible paths. The Android-ready linear
model is exported as a JSON artifact, while the one-class path exports a TFLite autoencoder
\citep{tensorflow_lite_guide}. These models are intentionally lightweight, but they serve mainly as
feasibility baselines relative to the stronger remote models rather than as the final thesis
headline.

Alert generation includes severity stabilization and false-positive budget enforcement, then stores
triage-ready records with explanation text and feature contributors.

\subsection{UI and Triage Workflow}
The Compose UI exposes:
\begin{itemize}[leftmargin=*]
  \item consent/disclosure and capture controls,
  \item backend URL/token and export toggle,
  \item alert center with triage states (\texttt{OPEN}, \texttt{INVESTIGATING},
  \texttt{RESOLVED}, \texttt{FALSE\_POSITIVE}),
  \item local data purge and anonymized CSV snapshot export.
\end{itemize}
This supports both standalone local operation and analyst-assisted workflows.

\section{Backend Adapter Implementation}
The backend adapter is a FastAPI service with token authentication, strict request validation, and
defense-in-depth response headers. Incoming events are enqueued in SQLite and forwarded with bounded
retry/backoff; exhausted items move to a dead-letter queue for replay.

Beyond ingest, the adapter implements:
\begin{itemize}[leftmargin=*]
  \item alert triage APIs and incident summaries,
  \item per-device policy storage and auto-tuning hooks,
  \item policy simulation and quality summary endpoints,
  \item retraining sample export and forensics bundle export,
  \item optional Wazuh/Resistine connector routes \citep{wazuh_api}.
\end{itemize}
This architecture isolates endpoint concerns from SIEM-specific integration complexity.

\section{ML Pipeline Implementation}
The ML pipeline is organized as reproducible command-line stages that now operate over a merged
public corpus rather than only over a controlled synthetic trace:
\begin{itemize}[leftmargin=*]
  \item public-dataset normalization and merged-corpus manifest generation,
  \item grouped split audits and evaluation-protocol reports,
  \item on-device linear and TFLite export,
  \item remote backend training across multiple model families,
  \item privacy ablation and metadata leakage attack benchmarking,
  \item encrypted-flow sequence fingerprint benchmarking,
  \item privacy-student training and federated simulation,
  \item comparison-matrix and performance-gate generation.
\end{itemize}

\subsection{Privacy Views}
The privacy mechanism is implemented as explicit feature views rather than as an afterthought.
\texttt{off} and \texttt{low} retain the full feature set, while \texttt{medium} and \texttt{strict}
replace many exact values with coarse or distribution-aware buckets. For example, activity,
volume, duration, rate, burstiness, stability, and shape are represented through a small number of
bucketized values, while direct host hints and several stronger destination/timing cues are removed
from the exported representation.

This is not packet padding or transport obfuscation. The system does not add dummy traffic. The
privacy transformation happens at the feature-representation level.

\subsection{Privacy Student and Federated Path}
The privacy-student model is trained on the reduced \texttt{medium} view with reconstruction-style
pretraining and multiple adversary heads. Those adversaries try to recover app identity, app
family, dataset source, context bucket, and destination behavior from the learned latent
representation. The objective is to preserve anomaly utility while reducing leakage of sensitive
traffic attributes.

The federated simulator then trains over the privacy-student latent representation using a FedProx
style algorithm rather than over raw full-feature windows. This does not prove secure deployment by
itself, but it makes the collaborative-learning design concrete and measurable.

\subsection{Reproducibility Support}
The \texttt{run\_experiment\_suite} path, cached privacy-view derivation, manifest reports, and
full comparison matrix are all designed to support reproducibility. The archived report
set used in this manuscript was produced from the March 20--21, 2026 public-corpus run and is
archived in the evidence bundle described in the appendix.
